\documentclass[12pt]{article}  
\newcommand{\say}[1]{``#1''}  
\newcommand{\nsay}[1]{`#1'}  
\usepackage{endnotes}  
\newcommand{\B}{\backslash{}}  
\renewcommand{\,}{\textsuperscript{,}}  
\usepackage{setspace}   
\usepackage{tipa}  
\usepackage{hyperref}  
\begin{document}  
\doublespacing  
\section{\href{feeling-heavy.html}{On Feeling Heavy}}  
First Published: 2026 April 5

\section{Draft 2: 5 April 2026}

I do think that there could be something interesting in my trying to see how tight the correlation is between the number of words in a draft or post that I take to get to its subject and the amount that I don't want to interact with the subject.  
I right now assume that the correlation would be very strong, only limited by the fact that I do not always have the same writing pace.

The first draft touches much more on the title of this post, which is feeling heavy.  
Right now, I'm remembering the second half of what I must do every time that I feel like this, or in this realm, at least: I need to ignore that voice.

Right now, the albatross that feels like an anchor is the fact that I have not been writing.  
Another albatross is that my home is not clean enough to host those I wish to host.

Those are both real and true things!  
It is not just valid that I feel them, it is if anything probably a good thing that I feel dissatisfaction with the state of my home and the absence of words I have written.\footnote{typed}  
However, spirals, fascinating as they are to look at when drawn, are just as captivating when dragging me down.

This month starts a biannual writing competition on the site I write\footnote{past and future tenses, so it averages to present, right?} my web serial.  
I told myself that I would use that to make me write again, and this feels like the reasoning I need right now.  
So, what are the fiction ideas I've told myself I would like to write some day?  
I guess you'll have to see!

\section{Draft 1: 5 April 2026}

So even though it has been far too long since I wrote anything, let alone on this exact site, I do want the record to show that I did write a full\footnote{more than one draft, 3400 or so words} post yesterday.  
But, I didn't feel like I had gotten to a stable point in my thoughts.   
I'm currently removed enough from the daily practice of this site to remember if I used to feel this sense of closure, but I'm going to imagine that I did.  
Even when I ended somewhere that wasn't a final answer, it at least felt like a final enough answer.  
Yesterday's post, perhaps because the first 2500 words were me attempting to discover what the post was going to be about, did not feel finalized.

My past writing is, in one sense, the topic of today's proflection\footnote{Is that the word I ended up settling on?}.  
In general, however, my goal for today's post is to externalize the internal feeling that I have right now.

In previous posts, I know that I've expressed, whether implicitly or explicitly\footnote{oof I'm loving parentheticals, or whatever else you'd call the sentence structure I'm using here}, the I have felt as though my body and mind betray me.  
I still somewhat deal with this struggle, though I am actively putting in the work day by day to accept that my body does not fail me.  
Either I failed it by failing to give it the nourishment it needed, or it was not a task that my body was capable of.\footnote{I know I've had the conversation about how failure is not a moral judgement, and I do stand by that, and that it's ok to fail, with at least one of my regular readers (ooooof that sentence structure hurts me but I've decided no more pressing the delete key unless it was literally a typo for now). However, I do think that it's not fair to call it a failure for a child to not stop a speeding train (oof don't make the mental image, just know I meant \say{a not Superman to do a Superman thing}. If I don't give someone any colors, just black ink, they did not fail to paint an accurate picture of something rosy red. They were not given what they need for that. So too my body (is the idea)}

N.B. to future me\footnote{rhyme!}: maybe try starting the inevitable draft two with a meta reflection on how the amount of time it takes me (or number of words, at least) to actually start talking about a thing here is probably correlated to how uncomfortable I am with it.

It's relatively easy for me to accept that my body has limitations, because I can see the effects of aging on myself, and I know that much of the athleticism I once was capable of required double digits of time spent actively working my body.  
That is not my priority right now, and so it's unfair to expect same results with less resource allocation.

It's much harder, however, for me to accept the ways it feels as though my mind fails me.

My memory feels terrible these days, and I don't like that.  
I can still remember facts, at least in the eyes of most, but my visual memory is effectively gone.  
I have very few moments I can see in my mind's eye, and those that I  can tend to be from a third-person perspective, letting me know that they're more imagination than recollection.  
I assume that this, like the fact that I've gotten worse at mental math, has more to do with where I focus my mental energy than much else, and so I can somewhat accept it.

But, last night I felt a migraine setting in.  
I, in what feels like a big change for me, left where I was and went to sleep, rather than trying to power through the pain and keep going, like I would have for much of my life.  
I don't like that my body needed the break there, especially because I cannot point to a cause.

But, more than that, as the title suggests, I have difficulty with the way that my mind makes everything feel heavy right now.  
Depending on the metaphor I've used, depression has been a variety of things.  
However, at least as far as I can remember, the load has felt metaphorical.  
These past few days and today in particular, I feel as though I am actively being pushed down.

Each time that I try to stand, there is a voice in my head that tells me it's pointless.  
There's another that points out that I'm already feeling kind of light headed, so do I really want to risk standing up?  
Another mentions that we've just gotten to a good point in whatever activity I was exercising.

To be clear, none of these voices are verbalizing within my mind.\footnote{Not in the \say{I promise I don't hear voices} way, but in the \say{these are subvocalizations I can track into full thoughts if I put forth the effort}}  
I have these vague impressions that I should not stand up.  
Vague may be the wrong word.  
I have amorphous but powerful impressions telling me to not move.

As the series of posts I have made about feeding myself might imply, I struggle to nourish my body.

Today in particular, when each breath actually feels like effort, when even the act of blinking seems to take a bit of effort, food is much more difficult.

But, more than anything, it's this weight that presses me down and tells me not only that I cannot do anything right now, but that anything I did would be pointless anyways.  
I have an unwelcome visitor right now who tells me that I have failed to live up to my own potential, that I have failed those who love me, that I have failed myself.  
I don't like feeling this way.

Right now, I think that it may be most acute, because, like most procrastinators, I had an external point where I told myself I would start writing again.  
The website where I used to write my web serial has a biannual writing competition, where authors attempt to put out a little over half a century thousand of words in thirty five days.  
It's something I've done before, while actively keeping up with my other novel.  
I think that my idea was that, realistically speaking, fifteen hundred words a day is not that many for me.\footnote{I'm at over a thousand right now, and the document claims I haven't even been on this page for 19 minutes yet}  
A shock to the system might be what I need to get back into this and the other writing I do, was the thought at least.

Of course, I did nothing to keep track of when, exactly, the event started.  
I realized yesterday, after writing the three thousand and some words about where I am right now, emotionally, that it had started on the first.  
Missing five days isn't really that much time to miss.\footnote{again, I type fast and the expectations are low for writing quality, both in myself and in the general reader base}

Honestly, part of what I always tell myself is that the weight I feel is not a sign that I should give up. It's a sign that I should ignore the pain and power through.  
Let's take that advice to heart, at least for now, and try to write these 55000 words before the next four and a touch weeks have gone by.  
I guess I need to figure out what I'd write about; I think I want to at least start a new story, even if I'm not totally sure where/what or anything else.  
That's something I can consider.  
Or, I can just start writing and see what happens?  
I can look through my notes and see what fiction ideas I've had in the past.

\end{document}